\begin{prob}

    The table below shows the number of nodes in a graph $G$ which are distance 0
    from node $s$, distance 1 from node $s$, distance 2, etc.:

    \begin{center} 
        \begin{tabular}{c|cccccc}\toprule
            Distance from $s$ & 0 & 1 & 2 & 3 & 4 & 5 \\ \midrule
            Number of nodes   & 1 & 7 & 3 & 6 & 3 & 1 \\ \bottomrule
        \end{tabular}
    \end{center}

    Suppose a BFS is performed on this graph with node $s$ as the starting node.
    What is the largest number of nodes that could possibly be in the queue at any one time?

    \inlineresponsebox[.75in]{9}

    \begin{soln}
        Remember that the queue can only simultaneously contain nodes at two different
        distances (it can't contain nodes at distance 1, distance 2, and distance 3
        at the same time, but it can contain nodes at distance 1 and distance 2 at the
        same time).

        The largest number of nodes that could possibly be in the queue at any one time
        is when we have (almost) all of the nodes from the first distance, and all
        of the nodes from the second distance.

        In this graph, let $u$ be a neighbor of the source, and suppose it has
        an edge with all nodes that are distance 2 from the source. If $u$ is
        the first node popped from the queue (after the source), then once its
        neighbors are added to the queue, the queue will contain all of the
        nodes at distance 1 except for $u$ (for 6 nodes in total), and all of
        the nodes at distance 2 (for 3 nodes in total). This gives us a total
        of 9 nodes in the queue.
    \end{soln}

\end{prob}
